\documentclass[11pt,a4paper]{article}

% ---------------------------------------------------------
% PREAMBLE (stable, warning-free, Overleaf-safe)
% ---------------------------------------------------------
\usepackage[utf8]{inputenc}
\usepackage[T1]{fontenc}
\usepackage{lmodern}
\usepackage{amsmath, amssymb, amsfonts}
\usepackage{bm}
\usepackage{geometry}
\usepackage{graphicx}
\usepackage{hyperref}
\usepackage{cite}
\usepackage{authblk}
\usepackage{physics}
\usepackage{mathtools}

\geometry{margin=2.4cm}

\title{\textbf{Companion XXXIV: Filaments, Coherence, and the Geometry of the Cosmic Web in the Relaxation-Driven Cyclic Cosmology}}
\author{Michael Lehmann \\ \texttt{mi.lehmann@gmx.de}}
\date{April 2026}

\begin{document}
\maketitle

\begin{abstract}
The cosmic web is the most striking manifestation of non-linear structure formation. 
Its filaments, sheets, and voids encode the dynamical history of the Universe and 
reflect the coherence of the underlying gravitational field. In the 
Relaxation-Driven Cyclic Cosmology (RDCC), the scale-dependent evolution of the 
gravitational potential alters the geometry and dynamics of the web in a 
characteristic way. This Companion develops a continuous description of filament 
formation, examines the influence of the relaxation scale on web coherence, and 
derives analytical expressions for filament thickness, curvature, and stability. 
The resulting picture reveals a cosmic web that is smoother, more coherent, and 
more strongly shaped by large-scale modes than in the standard cosmological model.
\end{abstract}
% ---------------------------------------------------------
\section{Introduction}

The cosmic web emerges from the anisotropic collapse of matter under gravity. 
Filaments form where collapse proceeds along two principal axes, sheets where only 
one axis collapses, and voids where expansion dominates. In the standard 
cosmological model, the geometry of the web is governed by the interplay between 
the initial tidal field and the non-linear amplification of density perturbations.

In RDCC, the relaxation-modified potential evolution introduces a new layer of 
structure. Large-scale modes retain a higher degree of temporal coherence, while 
small-scale modes experience a suppressed evolution. This imbalance reshapes the 
web's architecture: filaments become smoother, voids expand more coherently, and 
the connectivity of the web reflects the influence of the relaxation scale 
$k_{\mathrm{rel}}$.
% ---------------------------------------------------------
\section{Filament Formation in RDCC}

Filaments arise from the collapse of matter along two eigenvectors of the tidal 
tensor. The collapse time is determined by the evolution of the gravitational 
potential $\Phi(k,a)$, which in RDCC takes the form
\begin{equation}
    \Phi_{\mathrm{RDCC}}(k,a)
    = \Phi(k,a)
      \left[1 - \chi_{\Phi}
      \left(\frac{k}{k_{\mathrm{rel}}}\right)^{\alpha}\right].
\end{equation}

Modes with $k \ll k_{\mathrm{rel}}$ evolve coherently, while modes with 
$k \gg k_{\mathrm{rel}}$ experience a reduced potential decay. This difference 
modifies the collapse sequence. Filaments form earlier in regions dominated by 
large-scale coherence, and their cross-sectional profiles become smoother because 
the suppressed small-scale modes contribute less to the steepening of the density 
gradient.

The characteristic filament radius can be estimated by balancing the tidal 
compression with the relaxation-modified potential:
\begin{equation}
    r_{\mathrm{fil}}
    \sim \left[
        \frac{\Phi_{\mathrm{RDCC}}}{\lambda_{1} - \lambda_{2}}
    \right]^{1/2},
\end{equation}
where $\lambda_{1}$ and $\lambda_{2}$ are the dominant eigenvalues of the tidal 
tensor.
% ---------------------------------------------------------
\section{Coherence and Curvature of Filaments}

The coherence of a filament refers to the degree to which its spine follows a 
smooth trajectory through space. In RDCC, the enhanced temporal stability of 
large-scale modes leads to a more coherent gravitational field, which in turn 
reduces the curvature of filament spines.

The curvature $\kappa$ of a filament can be approximated by
\begin{equation}
    \kappa
    \sim \frac{\abs{\nabla_{\perp}\Phi_{\mathrm{RDCC}}}}{
        \abs{\nabla_{\parallel}\Phi_{\mathrm{RDCC}}}^{2}},
\end{equation}
where the subscripts denote derivatives perpendicular and parallel to the filament 
axis. Because the relaxation mechanism suppresses small-scale fluctuations, the 
perpendicular gradient becomes smoother, reducing $\kappa$ and producing straighter, 
more extended filaments.

This effect is particularly pronounced in regions where the tidal field is dominated 
by modes near $k_{\mathrm{rel}}$, which act as a stabilizing backbone for the web.
% ---------------------------------------------------------
\section{Void Expansion and Web Connectivity}

Voids expand as matter flows outward along the eigenvectors of the tidal tensor. 
In RDCC, the suppression of small-scale potential evolution reduces the formation 
of substructure within voids, allowing them to expand more coherently. The void 
radius evolves approximately as
\begin{equation}
    R_{\mathrm{void}}(a)
    \propto a^{\gamma_{\mathrm{RDCC}}},
\end{equation}
where $\gamma_{\mathrm{RDCC}}$ exceeds its standard-cosmology counterpart due to 
the enhanced influence of large-scale modes.

The connectivity of the web is also altered. Filaments linking massive halos become 
more stable, while the number of small-scale filaments decreases. This produces a 
web with fewer but more prominent connections, reflecting the imprint of the 
relaxation scale on the geometry of structure formation.
% ---------------------------------------------------------
\section{Observational Signatures}

The RDCC modifications to the cosmic web manifest in several measurable ways. 
Weak-lensing maps reveal smoother filamentary structures and a reduced abundance 
of small-scale peaks. Galaxy surveys show enhanced alignment of galaxies along 
filaments, reflecting the increased coherence of the underlying tidal field. 
Void statistics exhibit larger effective radii and a narrower distribution of 
shapes, consistent with the suppression of small-scale perturbations.

These signatures provide a direct observational window into the relaxation scale 
and the temporal structure of the gravitational field in RDCC.
% ---------------------------------------------------------
\section{Conclusion}

The cosmic web in the Relaxation-Driven Cyclic Cosmology is shaped by the 
scale-dependent evolution of the gravitational potential. Filaments form earlier 
and with smoother profiles, voids expand more coherently, and the connectivity of 
the web reflects the influence of the relaxation scale. These features distinguish 
RDCC from the standard cosmological model and offer a rich set of observational 
tests. The present Companion establishes the theoretical foundation for future 
studies of galaxy alignment, filament lensing, and the topology of the cosmic web 
within the RDCC framework.

% ========================
% References
% ========================

\begin{thebibliography}{10}

\bibitem{Flagship} Lehmann, M. (2026). Relaxation-Driven Cyclic Cosmology (RDCC): A Minimal CPT-Symmetric Framework with a Single Parameter. Flagship Paper. Zenodo: \url{https://zenodo.org/records/19744958} (v43.0 or later).

\bibitem{Zenodo} The complete RDCC companion series (I--LVII) is available at the same Zenodo record above.

% Thematisch relevante Companions
\bibitem{CompXVI} Companion XVI: Boltzmann Code and Modified Hierarchy.
\bibitem{CompXXXI} Companion XXXI: RDCC Halo Model and Non-Linear Structure.
\bibitem{CompXXXII} Companion XXXII--XXXIX: Galaxy--Halo Connection, Cosmic Web, Lensing, RSD, ISW, tSZ, Rotation Curves, CGM.

% Wichtige externe Referenzen
\bibitem{Planck2020} Planck Collaboration (2020), Astron. Astrophys. 641, A6.
\bibitem{DESI2024} DESI Collaboration (2024/2025), arXiv: [aktuelle $f\sigma_8$/BAO-Paper].
\bibitem{JWST2024} Maiolino et al. (2024), Nature (JWST high-$z$ galaxies/quasars).
\bibitem{Kaiser1987} Kaiser, N. (1987), Mon. Not. R. Astron. Soc. 227, 1 (RSD).
\bibitem{Bartelmann2001} Bartelmann, M. \& Schneider, P. (2001), Phys. Rep. 340, 291 (Weak Lensing Review).
\bibitem{HuOkamoto2002} Hu, W. \& Okamoto, T. (2002), Astrophys. J. 574, 566 (CMB Lensing).
\bibitem{LewisChallinor2006} Lewis, A. \& Challinor, A. (2006), Phys. Rep. 429, 1 (CMB Lensing Review).
\bibitem{NFW1997} Navarro, J. F., Frenk, C. S. \& White, S. D. M. (1997), Astrophys. J. 490, 493 (Halo Profiles).

\end{thebibliography}

\end{document}
